\newpage
\addcontentsline{toc}{chapter}{Quellen}
\renewcommand{\refname}{Quellen}   
\renewcommand{\bibname}{Quellen} % Lässt die eigene Überschrift des Literaturverzeichnisses verschwinden. Es wird nun "Quellen" angezeigt.
\bibliography{include/Quellen.bib}
\bibliographystyle{alphadin}

% Dieser Befehl kann genutzt werden um hinter einem Text einen Verweis auf die Quelle in der Quellenangabe erscheinen zu lassen.
%\cite{Example1}
% Dieser Befehl kann genutzt werden um Quellenangaben in den Quellen aufzulisten, auch wenn im Text nicht darauf verwiesen wird.
%\nocite{*}
\nocite{Vorbereitungshilfe}
\nocite{Ordner}
\section*{Abbildungsverzeichnis}
\autoref{fig:Aufbau}:  \cite{Vorbereitungshilfe}, p.199\\
\autoref{fig:Blockschaltbild}:  \cite{Vorbereitungshilfe}, p.201\\
\autoref{fig:decay_co60}:  \url{https://upload.wikimedia.org/wikipedia/commons/thumb/0/03/Cobalt-60m-decay.svg/1067px-Cobalt-60m-decay.svg.png}


